\documentclass[11pt]{article}

\usepackage[margin=1in]{geometry}
\usepackage[T1]{fontenc}
\usepackage[utf8]{inputenc}
\usepackage{amsmath,amssymb,amsthm,mathtools}
\usepackage{booktabs,array,longtable}
\usepackage{enumitem}
\usepackage{hyperref}
\usepackage{xcolor}
\usepackage{microtype}

\hypersetup{
  colorlinks=true,
  linkcolor=blue!55!black,
  citecolor=blue!55!black,
  urlcolor=blue!55!black,
  pdftitle={Computing the Cohomology Ring of the Infinite Grassmannian},
  pdfauthor={W. Huang, revised seminar note}
}

\setlength{\parindent}{0pt}
\setlength{\parskip}{0.65em}
\setlist[itemize]{topsep=0.25em,itemsep=0.2em}
\setlist[enumerate]{topsep=0.25em,itemsep=0.2em}

\newtheorem{theorem}{Theorem}[section]
\newtheorem{proposition}[theorem]{Proposition}
\newtheorem{lemma}[theorem]{Lemma}
\newtheorem{corollary}[theorem]{Corollary}
\theoremstyle{remark}
\newtheorem{remark}[theorem]{Remark}
\newtheorem{example}[theorem]{Example}

\newcommand{\F}{\mathbb F}
\newcommand{\R}{\mathbb R}
\newcommand{\RP}{\mathbb{RP}}
\newcommand{\Gr}{\operatorname{Gr}}
\newcommand{\Gn}{G_n}
\newcommand{\gamman}{\gamma_n}
\newcommand{\lambdainf}{\lambda^\infty}
\newcommand{\lambdaone}{\lambda^1}
\newcommand{\wt}{\widetilde{w}}
\newcommand{\olW}{\overline{W}}
\newcommand{\parts}{\operatorname{part}}

\newenvironment{talknote}{\begin{quote}\small\textbf{Speaker note.} }{\end{quote}}
\newenvironment{optional}{\begin{quote}\small\textbf{Optional material.} }{\end{quote}}

\title{Computing $H^*(\Gr_n(\R^\infty);\F_2)$ and Uniqueness of Stiefel--Whitney Classes\\[0.4em]
\large Revised 90-minute seminar talk note}
\author{W. Huang\\Seminar: Characteristic Classes}
\date{12 June 2026}

\begin{document}
\maketitle

\begin{abstract}
This note is written for a graduate seminar talk whose assigned topic is to compute
\[
  H^*(\Gr_n(\R^\infty);\F_2)
\]
and to prove uniqueness of the Stiefel--Whitney classes, following the standard approach in Milnor--Stasheff, Section 7.  We write
\[
  \F_2=\mathbb Z/2\mathbb Z,\qquad \Gn=\Gr_n(\R^\infty),
\]
and let
\[
  \gamman\to \Gn
\]
be the tautological rank $n$ real vector bundle.  The main result is
\[
  H^*(\Gn;\F_2)\cong \F_2[w_1,\dots,w_n],\qquad |w_i|=i,
\]
where $w_i=w_i(\gamman)$.

The proof has two main ingredients.  First, pull back to the split bundle over $(\RP^\infty)^n$ to show that $w_1,\dots,w_n$ satisfy no polynomial relation.  Second, use the Schubert cell structure of the infinite Grassmannian to show that these classes already account for all cohomology classes.  The same computation then gives a clean proof that the Stiefel--Whitney axioms characterize the classes uniquely.
\end{abstract}

\tableofcontents

\section*{Suggested 90-minute schedule}
\addcontentsline{toc}{section}{Suggested 90-minute schedule}

\begin{longtable}{@{}p{0.16\textwidth}p{0.78\textwidth}@{}}
\toprule
Time & Content \\
\midrule
0--5 min & State the goal: compute the mod $2$ cohomology ring of $\Gn=\Gr_n(\R^\infty)$, interpret it as the ring of mod $2$ characteristic classes of rank $n$ real vector bundles, and prove uniqueness of Stiefel--Whitney classes. \\
5--15 min & Recall the classification theorem for vector bundles, the tautological bundle $\gamman$, and why characteristic classes of rank $n$ bundles are the same as cohomology classes of $\Gn$. \\
15--25 min & List the inputs from earlier talks: the Stiefel--Whitney axioms, $H^*(\RP^\infty;\F_2)$, the Kunneth theorem over $\F_2$, the Schubert cell count for $\Gn$, and the fundamental theorem of symmetric polynomials. \\
25--45 min & Main proof, Part I: over $X=(\RP^\infty)^n$, construct the split bundle $\xi=\ell_1\oplus\cdots\oplus\ell_n$, compute $w(\xi)=\prod_i(1+a_i)$, and prove algebraic independence of $w_1(\gamman),\dots,w_n(\gamman)$. \\
45--60 min & Main proof, Part II: use Schubert cell counting to bound $\dim_{\F_2}H^r(\Gn;\F_2)$, compare with the number of monomials in the $w_i$'s of degree $r$, and conclude that $H^*(\Gn;\F_2)=\F_2[w_1,\dots,w_n]$. \\
60--68 min & Work out one or two low-rank examples, especially $n=1$ and $n=2$.  Treat $n=3$ as optional if time is short. \\
68--82 min & Prove uniqueness of Stiefel--Whitney classes from the ring computation. \\
82--90 min & Summarize the proof, answer questions, and emphasize that this talk proves uniqueness, while existence is a separate construction. \\
\bottomrule
\end{longtable}

\textbf{One-sentence goal.}  The infinite Grassmannian $\Gn$ is the classifying space for rank $n$ real vector bundles, and its mod $2$ cohomology ring is exactly the polynomial ring generated by the universal Stiefel--Whitney classes.

\begin{talknote}
The two most important phrases to repeat during the talk are: ``test on a split bundle'' and ``compare dimensions using Schubert cells.''  These are the two ideas that hold the proof together.
\end{talknote}

\section{Motivation and context}

A vector bundle may be locally trivial but globally twisted.  Characteristic classes are cohomology classes that detect part of this global twisting.  For a real vector bundle $\xi\to B$, the Stiefel--Whitney classes are classes
\[
  w_i(\xi)\in H^i(B;\F_2).
\]
They are natural in the base space and behave multiplicatively under direct sum.  The total Stiefel--Whitney class is
\[
  w(\xi)=1+w_1(\xi)+w_2(\xi)+\cdots .
\]

The point of this talk is not to construct Stiefel--Whitney classes.  We assume that classes satisfying the usual axioms exist.  Under that assumption, we compute the universal ring in which all rank $n$ Stiefel--Whitney classes live.  Then we prove that the axioms determine the classes uniquely.  Thus the logical structure is
\[
  \text{axioms plus existence}\quad\Longrightarrow\quad
  \text{universal computation}\quad\Longrightarrow\quad
  \text{uniqueness}.
\]
The actual construction of the classes is a separate topic.

\begin{talknote}
A good opening is: ``In the previous talk, we studied Schubert cells in real Grassmannians.  Today we use that cell decomposition to compute a cohomology ring.  The answer is surprisingly simple: the mod $2$ cohomology of the infinite Grassmannian is a polynomial ring, and the generators are the universal Stiefel--Whitney classes.''
\end{talknote}

\section{Notation and standing facts}

Throughout the talk, set
\[
  \F_2=\mathbb Z/2\mathbb Z,
  \qquad
  \Gn=\Gr_n(\R^\infty).
\]
The space $\Gn$ is the infinite real Grassmannian of $n$-planes in $\R^\infty$.  It carries the tautological rank $n$ vector bundle
\[
  \gamman\longrightarrow \Gn,
\]
whose fiber over an $n$-plane $V\subset \R^\infty$ is $V$ itself.

We write
\[
  w_i=w_i(\gamman)\in H^i(\Gn;\F_2)
\]
when there is no risk of confusion.  Thus the object to compute is the graded ring $H^*(\Gn;\F_2)$, and the expected answer is
\[
  H^*(\Gn;\F_2)\cong \F_2[w_1,\dots,w_n],\qquad |w_i|=i.
\]

\subsection{Classification of vector bundles}

A basic input from the theory of classifying spaces is the following theorem.

\begin{theorem}[Classification of real vector bundles]
Let $B$ be a paracompact space.  Isomorphism classes of rank $n$ real vector bundles over $B$ are in bijection with homotopy classes of maps
\[
  [B,\Gn].
\]
More precisely, if $\xi\to B$ is a rank $n$ real vector bundle, then there is a classifying map
\[
  f:B\to \Gn
\]
such that
\[
  \xi\cong f^*\gamman.
\]
The map $f$ is unique up to homotopy.
\end{theorem}

This theorem explains why $\Gn$ is the correct universal object.  A cohomology class
\[
  c\in H^k(\Gn;\F_2)
\]
defines a characteristic class for rank $n$ bundles by the rule
\[
  c(\xi)=f^*c\in H^k(B;\F_2),
\]
where $f:B\to \Gn$ classifies $\xi$.  Homotopy invariance makes this well-defined.

Conversely, every natural characteristic class for rank $n$ bundles is determined by its value on the universal bundle $\gamman$.  Therefore the ring of mod $2$ characteristic classes for rank $n$ real vector bundles is naturally identified with
\[
  H^*(\Gn;\F_2).
\]
So computing $H^*(\Gn;\F_2)$ is the same thing as computing all mod $2$ characteristic classes of rank $n$ real vector bundles.

\subsection{Axioms for Stiefel--Whitney classes}

We use the following properties of Stiefel--Whitney classes.

\begin{enumerate}
\item \textbf{Degree and rank condition.}  If $\xi\to B$ has rank $r$, then
\[
  w_i(\xi)\in H^i(B;\F_2),\qquad w_0(\xi)=1,
  \qquad w_i(\xi)=0\text{ for }i>r.
\]

\item \textbf{Naturality.}  If $f:B'\to B$, then
\[
  w_i(f^*\xi)=f^*w_i(\xi).
\]

\item \textbf{Whitney product formula.}  If $\xi$ and $\eta$ are real vector bundles over the same base, then
\[
  w(\xi\oplus\eta)=w(\xi)w(\eta).
\]

\item \textbf{Line-bundle normalization.}  If $\lambda^1\to \RP^1$ is the tautological real line bundle and $u\in H^1(\RP^1;\F_2)$ is the nonzero element, then
\[
  w(\lambda^1)=1+u.
\]
\end{enumerate}

Because we work over $\F_2$, there are no sign issues in cup products.  We usually omit the cup symbol and write products multiplicatively.

\subsection{Cohomology of $\RP^\infty$}

We use the standard computation
\[
  H^*(\RP^\infty;\F_2)\cong \F_2[a],\qquad |a|=1.
\]
Here $a=w_1(\lambdainf)$, where $\lambdainf\to \RP^\infty$ is the tautological line bundle.

For the product
\[
  X=(\RP^\infty)^n,
\]
let
\[
  \pi_i:X\to \RP^\infty
\]
be the $i$-th projection and define
\[
  a_i=\pi_i^*a\in H^1(X;\F_2).
\]
The Kunneth theorem over the field $\F_2$ gives
\[
  H^*(X;\F_2)\cong \F_2[a_1,\dots,a_n].
\]
This polynomial ring will be our testing ground for the universal classes $w_i$.

\section{Statement of the main theorem}

\begin{theorem}[Cohomology of the infinite Grassmannian]\label{thm:main}
For $\Gn=\Gr_n(\R^\infty)$, the Stiefel--Whitney classes of the tautological bundle give an isomorphism of graded rings
\[
  \F_2[x_1,\dots,x_n]\xrightarrow{\ \cong\ }H^*(\Gn;\F_2),
  \qquad x_i\longmapsto w_i(\gamman),
\]
where $|x_i|=i$.  Equivalently,
\[
  H^*(\Gn;\F_2)=\F_2[w_1,
  \dots,w_n],\qquad |w_i|=i.
\]
\end{theorem}

The proof has two parts.

First, we show that $w_1,\dots,w_n$ satisfy no polynomial relation.  This says that the map
\[
  \F_2[x_1,\dots,x_n]\longrightarrow H^*(\Gn;\F_2),
  \qquad x_i\longmapsto w_i(\gamman),
\]
is injective.

Second, we show that there is no room for any additional cohomology classes.  This uses the Schubert cell decomposition of $\Gn$.  In degree $r$, the number of possible monomials in the $w_i$'s is exactly the number of Schubert cells of dimension $r$.  Thus the injective polynomial subring already has the maximum possible dimension in each degree.

\section{Part I: algebraic independence}

\subsection{The split bundle over $(\RP^\infty)^n$}

Let
\[
  X=(\RP^\infty)^n.
\]
For each $i$, let
\[
  \ell_i=\pi_i^*\lambdainf
\]
be the pullback of the tautological line bundle from the $i$-th factor.  Now form the direct sum
\[
  \xi=\ell_1\oplus\ell_2\oplus\cdots\oplus\ell_n.
\]
This is a rank $n$ real vector bundle over $X$, split as a sum of line bundles.

By naturality,
\[
  w(\ell_i)=w(\pi_i^*\lambdainf)=\pi_i^*w(\lambdainf)=1+a_i.
\]
By the Whitney product formula,
\[
  w(\xi)=w(\ell_1)\cdots w(\ell_n)=\prod_{i=1}^n(1+a_i).
\]
Expanding this product gives
\[
  w(\xi)=1+e_1(a_1,\dots,a_n)+e_2(a_1,\dots,a_n)+\cdots+e_n(a_1,\dots,a_n),
\]
where $e_k$ is the $k$-th elementary symmetric polynomial:
\[
  e_k(a_1,\dots,a_n)=\sum_{1\le i_1<\cdots<i_k\le n}a_{i_1}\cdots a_{i_k}.
\]
Therefore
\[
  w_k(\xi)=e_k(a_1,\dots,a_n).
\]

\begin{example}
For $n=3$, the computation says
\[
  w(\ell_1\oplus\ell_2\oplus\ell_3)=(1+a_1)(1+a_2)(1+a_3),
\]
so
\[
  w_1=a_1+a_2+a_3,
\]
\[
  w_2=a_1a_2+a_1a_3+a_2a_3,
\]
\[
  w_3=a_1a_2a_3.
\]
Thus the universal Stiefel--Whitney classes will pull back to the elementary symmetric polynomials.
\end{example}

\subsection{Pulling back the universal bundle}

Since $\xi\to X$ is a rank $n$ bundle, it has a classifying map
\[
  g:X\to \Gn
\]
such that
\[
  \xi\cong g^*\gamman.
\]
By naturality,
\[
  g^*w_k(\gamman)=w_k(g^*\gamman)=w_k(\xi)=e_k(a_1,\dots,a_n).
\]
Thus the induced map
\[
  g^*:H^*(\Gn;\F_2)\to H^*(X;\F_2)=\F_2[a_1,\dots,a_n]
\]
sends the universal classes to elementary symmetric polynomials:
\[
  g^*(w_k)=e_k(a_1,\dots,a_n).
\]

This is the key computational bridge.  We do not yet know all of $H^*(\Gn;\F_2)$, and at this point we are not claiming that $g^*$ is injective on all cohomology.  We only know exactly what happens to the particular classes $w_k$ after pulling them back to $X$.

\subsection{Elementary symmetric polynomials over $\F_2$}

We use the following algebra fact.

\begin{lemma}\label{lem:elementary-independent}
The elementary symmetric polynomials
\[
  e_1(a_1,\dots,a_n),\dots,e_n(a_1,\dots,a_n)
\]
are algebraically independent in $\F_2[a_1,\dots,a_n]$.
\end{lemma}

\begin{proof}
The fundamental theorem of symmetric polynomials applies over every commutative coefficient ring, in particular over $\F_2$.  It says that the map
\[
  \F_2[t_1,\dots,t_n]\longrightarrow \F_2[a_1,\dots,a_n]^{S_n},
  \qquad t_i\longmapsto e_i(a_1,\dots,a_n),
\]
is an isomorphism.  Equivalently, every symmetric polynomial in $a_1,\dots,a_n$ can be written uniquely as a polynomial in $e_1,\dots,e_n$.

The uniqueness part means precisely that there is no nonzero polynomial $P\in \F_2[x_1,\dots,x_n]$ such that
\[
  P(e_1,\dots,e_n)=0.
\]
Hence $e_1,\dots,e_n$ are algebraically independent.
\end{proof}

\begin{remark}
No division by $n!$ or by $|S_n|$ is involved.  This is why the theorem remains valid in characteristic $2$.
\end{remark}

\begin{talknote}
This is a common place for questions.  Say explicitly: ``The fundamental theorem of symmetric polynomials is an integral statement.  We are not averaging over the symmetric group, so characteristic $2$ causes no problem.''
\end{talknote}

\subsection{Injectivity of the polynomial map}

Define a graded ring homomorphism
\[
  \Phi:\F_2[x_1,\dots,x_n]\to H^*(\Gn;\F_2),
  \qquad x_i\longmapsto w_i(\gamman).
\]
We prove that $\Phi$ is injective.

Suppose
\[
  P(w_1,\dots,w_n)=0
\]
for some polynomial $P\in\F_2[x_1,\dots,x_n]$.  Pull back by $g:X\to \Gn$.  Then
\[
  0=g^*P(w_1,\dots,w_n)
   =P(g^*w_1,\dots,g^*w_n)
   =P(e_1,\dots,e_n)
\]
in $\F_2[a_1,\dots,a_n]$.  Since $e_1,\dots,e_n$ are algebraically independent, it follows that $P=0$.

Therefore there is no nontrivial polynomial relation among the universal classes $w_1,\dots,w_n$.

\begin{proposition}\label{prop:poly-injective}
The map
\[
  \F_2[x_1,\dots,x_n]\to H^*(\Gn;\F_2),
  \qquad x_i\longmapsto w_i(\gamman),
\]
is injective.
\end{proposition}

\begin{talknote}
Explain the point verbally as follows.  We tested the universal classes on a bundle that is as split as possible: a direct sum of tautological line bundles.  On this test bundle, the total Stiefel--Whitney class is literally the product $\prod_i(1+a_i)$.  Thus the universal classes become elementary symmetric polynomials.  Since those are algebraically independent, the original universal classes must also be algebraically independent.
\end{talknote}

\section{Part II: the cell-counting argument}

Algebraic independence gives us a polynomial subring
\[
  \F_2[w_1,\dots,w_n]\subseteq H^*(\Gn;\F_2).
\]
The next step is to prove that this subring is the whole cohomology ring.

\subsection{Schubert cells in $\Gn$}

From the Schubert cell decomposition of the infinite Grassmannian, the cells of $\Gn$ are indexed by partitions
\[
  \lambda=(\lambda_1,\dots,\lambda_n)
\]
with
\[
  \lambda_1\ge \lambda_2\ge\cdots\ge \lambda_n\ge 0.
\]
The dimension of the cell indexed by $\lambda$ is
\[
  |\lambda|=\lambda_1+\cdots+\lambda_n.
\]
Thus the number of cells of dimension $r$ is the number of partitions of $r$ into at most $n$ parts.  Denote this number by
\[
  p_{\le n}(r).
\]
There are only finitely many such cells in each fixed dimension $r$.

Cellular cochains give an immediate upper bound:
\[
  \dim_{\F_2}H^r(\Gn;\F_2)
  \le \dim_{\F_2}C^r(\Gn;\F_2)
  = \#\{\text{$r$-cells}\}=p_{\le n}(r).
\]
We do not need to compute the cellular differential.  The dimension of cohomology is always at most the dimension of the cochain group.

\subsection{Counting monomials in $w_1,\dots,w_n$}

Because $|w_i|=i$, a monomial
\[
  w_1^{r_1}w_2^{r_2}\cdots w_n^{r_n}
\]
has total degree
\[
  r_1+2r_2+\cdots+nr_n.
\]
Therefore the number of monomials in $w_1,\dots,w_n$ of total degree $r$ is the number of solutions
\[
  r_1+2r_2+\cdots+nr_n=r,
  \qquad r_i\ge 0.
\]
This is the number of partitions of $r$ whose parts have size at most $n$.

There is a standard bijection between
\[
  \{\text{partitions of $r$ into parts of size at most $n$}\}
\]
and
\[
  \{\text{partitions of $r$ into at most $n$ parts}\}.
\]
The bijection is conjugation of Young diagrams.  Hence the number of degree $r$ monomials in the $w_i$'s is also
\[
  p_{\le n}(r).
\]

\begin{example}
Let $n=2$ and $r=4$.  The degree $4$ monomials in $\F_2[w_1,w_2]$ are
\[
  w_1^4,
  \qquad w_1^2w_2,
  \qquad w_2^2.
\]
There are three of them.  The partitions of $4$ into at most two parts are
\[
  4,
  \qquad 3+1,
  \qquad 2+2.
\]
Again there are three.  This is the equality of counts used in the proof.
\end{example}

\begin{remark}
The argument compares numbers.  It does not identify Schubert classes with monomials in the $w_i$'s.  The proof only needs the upper bound from cellular cochains and the lower bound from algebraic independence.
\end{remark}

\subsection{Finishing the proof of the main theorem}

Since the map
\[
  \F_2[w_1,\dots,w_n]\hookrightarrow H^*(\Gn;\F_2)
\]
is injective, the degree $r$ monomials in the $w_i$'s are linearly independent in $H^r(\Gn;\F_2)$.  There are $p_{\le n}(r)$ such monomials.  Hence
\[
  \dim_{\F_2}H^r(\Gn;\F_2)
  \ge p_{\le n}(r).
\]
On the other hand, the cellular bound gives
\[
  \dim_{\F_2}H^r(\Gn;\F_2)
  \le p_{\le n}(r).
\]
Therefore
\[
  \dim_{\F_2}H^r(\Gn;\F_2)=p_{\le n}(r)
\]
for every $r$, and the degree $r$ monomials in $w_1,\dots,w_n$ form a basis for $H^r(\Gn;\F_2)$.

Thus the polynomial subring generated by $w_1,\dots,w_n$ is all of $H^*(\Gn;\F_2)$.  This proves
\[
  H^*(\Gn;\F_2)\cong \F_2[w_1,\dots,w_n],\qquad |w_i|=i.
\]

\begin{talknote}
The most important board picture in this part is a Young diagram and its transpose.  Put the two counts side by side: Schubert cells count partitions with at most $n$ parts; monomials count partitions with largest part at most $n$; conjugation identifies them.
\end{talknote}

\section{Low-rank examples}

This section is useful for oral presentation, but it is not logically necessary for the proof.  If time is short, do only $n=1$ and $n=2$.

\subsection{The case $n=1$}

For $n=1$,
\[
  G_1=\Gr_1(\R^\infty)=\RP^\infty.
\]
The theorem says
\[
  H^*(G_1;\F_2)=\F_2[w_1].
\]
This is just the familiar computation
\[
  H^*(\RP^\infty;\F_2)=\F_2[a],
\]
with $w_1(\lambda^\infty)=a$.

This example is useful because it fixes the normalization of Stiefel--Whitney classes.  A real line bundle has only one possibly nonzero positive-degree Stiefel--Whitney class, namely $w_1$.

\subsection{The case $n=2$}

For $n=2$,
\[
  H^*(G_2;\F_2)=\F_2[w_1,w_2],
  \qquad |w_1|=1,
  \quad |w_2|=2.
\]
The first few graded pieces have bases
\[
\begin{array}{c|l}
\text{degree} & \text{basis}\\
\hline
0 & 1\\
1 & w_1\\
2 & w_1^2,\ w_2\\
3 & w_1^3,\ w_1w_2\\
4 & w_1^4,\ w_1^2w_2,\ w_2^2.
\end{array}
\]
This is a concrete illustration of the theorem: every cohomology class is a polynomial in the universal $w_1$ and $w_2$, and there are no relations.

\subsection{The case $n=3$}

For $n=3$,
\[
  H^*(G_3;\F_2)=\F_2[w_1,w_2,w_3].
\]
In degree $4$, a basis is
\[
  w_1^4,
  \qquad w_1^2w_2,
  \qquad w_2^2,
  \qquad w_1w_3.
\]
The four corresponding partitions of $4$ into at most three parts are
\[
  4,
  \qquad 3+1,
  \qquad 2+2,
  \qquad 2+1+1.
\]

\section{Interpretation as the ring of characteristic classes}

Let us spell out the meaning of the phrase ``the ring of $\mathbb Z/2\mathbb Z$-characteristic classes.''  Fix rank $n$.  A mod $2$ characteristic class of degree $k$ is a natural rule assigning to every rank $n$ real vector bundle $\xi\to B$ a class
\[
  c(\xi)\in H^k(B;\F_2).
\]
Naturality means that
\[
  c(f^*\xi)=f^*c(\xi)
\]
for every map $f:B'\to B$.

Given a class $c\in H^k(\Gn;\F_2)$, define
\[
  c(\xi)=f^*c,
\]
where $f:B\to \Gn$ classifies $\xi$.  This gives a characteristic class.  Conversely, a characteristic class is determined by its value on $\gamman$.  Thus, for paracompact base spaces,
\[
  \{\text{mod $2$ characteristic classes of rank $n$ real bundles}\}
  \cong H^*(\Gn;\F_2).
\]

Our computation says that every such characteristic class is a unique polynomial in
\[
  w_1,\dots,w_n.
\]
For example, for rank $3$ bundles, a degree $4$ characteristic class must have the form
\[
  \alpha w_1^4+\beta w_1^2w_2+\tau w_2^2+\sigma w_1w_3,
  \qquad \alpha,\beta,\tau,\sigma\in \F_2.
\]
No other degree $4$ characteristic classes exist.

\section{Uniqueness of Stiefel--Whitney classes}

We now use the ring computation to prove uniqueness.  The statement is that the Stiefel--Whitney axioms determine the classes completely.

\begin{theorem}[Uniqueness]\label{thm:uniqueness}
Suppose
\[
  W(\xi)=1+W_1(\xi)+W_2(\xi)+\cdots
\]
and
\[
  \olW(\xi)=1+\olW_1(\xi)+\olW_2(\xi)+\cdots
\]
are two assignments of total Stiefel--Whitney classes to real vector bundles, and suppose both assignments satisfy the Stiefel--Whitney axioms.  Then
\[
  W_i(\xi)=\olW_i(\xi)
\]
for every real vector bundle $\xi$ and every $i\ge 0$.
\end{theorem}

\subsection{Step 1: line bundles}

Let
\[
  \lambdainf\to \RP^\infty
\]
be the tautological line bundle over $\RP^\infty$, and let
\[
  \lambdaone\to \RP^1
\]
be the tautological line bundle over $\RP^1$.  Let
\[
  j:\RP^1\hookrightarrow \RP^\infty
\]
be the standard inclusion.  Then
\[
  j^*\lambdainf\cong \lambdaone.
\]

The normalization axiom over $\RP^1$ says that both assignments give
\[
  W(\lambdaone)=1+u,
  \qquad
  \olW(\lambdaone)=1+u,
\]
where $u$ is the nonzero element of $H^1(\RP^1;\F_2)$.  The map
\[
  j^*:H^1(\RP^\infty;\F_2)\to H^1(\RP^1;\F_2)
\]
is an isomorphism.  Therefore naturality forces
\[
  W(\lambdainf)=1+a,
  \qquad
  \olW(\lambdainf)=1+a.
\]
The rank condition also ensures that there are no higher Stiefel--Whitney classes for a line bundle.  Hence the two assignments agree on the tautological line bundle over $\RP^\infty$.

\subsection{Step 2: split bundles over $(\RP^\infty)^n$}

Return to
\[
  X=(\RP^\infty)^n,
  \qquad
  \xi=\ell_1\oplus\cdots\oplus\ell_n,
\]
where
\[
  \ell_i=\pi_i^*\lambdainf.
\]
Since the two assignments agree on $\lambdainf$, naturality implies that they agree on each $\ell_i$.  The Whitney product formula then gives
\[
  W(\xi)=\prod_{i=1}^n(1+a_i)=\olW(\xi).
\]
Thus the two assignments agree on this split test bundle.

\subsection{Step 3: pull back from the universal bundle}

Let
\[
  g:X\to \Gn
\]
classify $\xi$.  Then
\[
  \xi\cong g^*\gamman.
\]
By naturality,
\[
  g^*W(\gamman)=W(\xi)=\olW(\xi)=g^*\olW(\gamman).
\]
So
\[
  g^*\bigl(W(\gamman)-\olW(\gamman)\bigr)=0.
\]
Over $\F_2$, subtraction is the same as addition, but the notation is harmless.

Now use the ring computation.  Write
\[
  w_i=W_i(\gamman).
\]
Theorem~\ref{thm:main} gives
\[
  H^*(\Gn;\F_2)=\F_2[w_1,\dots,w_n],
\]
and the map $g^*$ sends
\[
  w_i\longmapsto e_i(a_1,\dots,a_n).
\]
Therefore $g^*$ is injective on $H^*(\Gn;\F_2)$.  Indeed, every class in $H^*(\Gn;\F_2)$ is uniquely $P(w_1,\dots,w_n)$, and if
\[
  g^*P(w_1,\dots,w_n)=P(e_1,\dots,e_n)=0,
\]
then algebraic independence of the elementary symmetric polynomials implies $P=0$.

Hence
\[
  W(\gamman)=\olW(\gamman).
\]

\subsection{Step 4: arbitrary bundles}

Let $\eta\to B$ be any rank $n$ real vector bundle, and let
\[
  f:B\to \Gn
\]
be a classifying map.  Then
\[
  \eta\cong f^*\gamman.
\]
Using naturality and equality on the universal bundle,
\[
  W(\eta)=f^*W(\gamman)=f^*\olW(\gamman)=\olW(\eta).
\]
Thus the two assignments agree on every rank $n$ bundle.  Since $n$ was arbitrary, they agree on all real vector bundles.

\begin{talknote}
For the uniqueness proof, keep the four steps very explicit: line bundles, split test bundle, universal bundle, arbitrary pullback.  The crucial sentence is: ``After the ring computation, $g^*$ is injective because every class is a polynomial in the $w_i$'s and $g^*w_i=e_i$.''
\end{talknote}

\section{Logical structure of the proof}

It may help to separate the proof into four conceptual moves.

\begin{enumerate}
\item \textbf{Universal viewpoint.}  Since $\Gn$ classifies rank $n$ bundles, the ring $H^*(\Gn;\F_2)$ is the ring of rank $n$ mod $2$ characteristic classes.

\item \textbf{Splitting test.}  Pull back $\gamman$ to $(\RP^\infty)^n$, where the test bundle is a sum of line bundles.  The total Stiefel--Whitney class becomes
\[
  \prod_{i=1}^n(1+a_i).
\]

\item \textbf{Algebraic independence.}  The classes $w_i$ pull back to elementary symmetric polynomials.  Since the elementary symmetric polynomials are algebraically independent, the $w_i$'s are algebraically independent.

\item \textbf{Dimension comparison.}  Schubert cells give an upper bound on $\dim H^r(\Gn;\F_2)$, and monomials in the $w_i$'s meet that upper bound exactly.
\end{enumerate}

Together these moves prove
\[
  H^*(\Gn;\F_2)=\F_2[w_1,\dots,w_n].
\]
The same splitting test and injectivity argument then prove uniqueness.

\section{Common questions and short answers}

\textbf{Why use $\F_2=\mathbb Z/2\mathbb Z$?}

Stiefel--Whitney classes are naturally mod $2$ characteristic classes.  Working over $\F_2$ also removes sign issues and makes the cohomology of $\RP^\infty$ equal to the simple polynomial ring $\F_2[a]$ with $|a|=1$.

\textbf{Why does cell counting give only an upper bound?}

The cellular cochain group $C^r(\Gn;\F_2)$ has one basis element for each $r$-cell, so $\dim C^r$ is the number of $r$-cells.  Cohomology is a quotient of cocycles, and cocycles are a subspace of cochains.  Hence
\[
  \dim H^r\le \dim C^r.
\]

\textbf{Do we need to know the cellular differential?}

No.  The injective polynomial subring already gives enough linearly independent cohomology classes to reach the cellular upper bound.  Therefore the dimension is forced, without computing the differential.

\textbf{Where is the previous talk used?}

It is used in the Schubert cell count: the number of $r$-dimensional cells in $\Gn$ is the number of partitions of $r$ into at most $n$ parts.

\textbf{Where are the Stiefel--Whitney axioms used?}

Naturality is used to pull classes back along maps and classifying maps.  The Whitney product formula is used to compute the split bundle over $(\RP^\infty)^n$.  The line-bundle normalization fixes the factor $1+a_i$.

\textbf{Does this prove existence?}

No.  The argument assumes classes satisfying the axioms exist, then proves the universal cohomology computation and uniqueness.  Existence is a separate construction.

\section{Optional details for questions}

\subsection{A slightly more detailed proof of the cell count}

A partition of $r$ into at most $n$ parts is a sequence
\[
  \lambda=(\lambda_1,\dots,\lambda_n)
\]
with
\[
  \lambda_1\ge\lambda_2\ge\cdots\ge\lambda_n\ge 0,
  \qquad
  \lambda_1+\cdots+\lambda_n=r.
\]
This is exactly the indexing set for the Schubert cells in $\Gn$, and $|\lambda|=r$ is the dimension of the corresponding cell.

On the other hand, a monomial in $\F_2[w_1,\dots,w_n]$ of degree $r$ has the form
\[
  w_1^{r_1}w_2^{r_2}\cdots w_n^{r_n}
\]
with
\[
  r_1+2r_2+\cdots+nr_n=r.
\]
This is the same as a partition of $r$ with $r_i$ parts equal to $i$.  Thus it is a partition of $r$ whose largest part is at most $n$.

Conjugating Young diagrams exchanges
\[
  \text{largest part at most $n$}
\]
with
\[
  \text{number of parts at most $n$}.
\]
Therefore the two counts agree.  This is the entire combinatorial content of the dimension comparison.

\subsection{Another view of algebraic independence}

Here is a direct way to think about why the elementary symmetric polynomials are algebraically independent.

Give $\F_2[a_1,
\dots,a_n]$ the lexicographic monomial order with
\[
  a_1>a_2>\cdots>a_n.
\]
Then the leading term of $e_k$ is
\[
  a_1a_2\cdots a_k.
\]
For a monomial
\[
  e_1^{m_1}e_2^{m_2}\cdots e_n^{m_n},
\]
the leading term is
\[
  a_1^{m_1+m_2+\cdots+m_n}
  a_2^{m_2+\cdots+m_n}
  \cdots
  a_n^{m_n}.
\]
The exponent vector uniquely determines $(m_1,\dots,m_n)$.  Thus distinct monomials in the $e_i$'s have distinct leading terms.  A nontrivial polynomial relation among the $e_i$'s could not cancel its largest leading term.  Hence no such relation exists.

This proof avoids quoting the full fundamental theorem of symmetric polynomials, though that theorem is the standard reference.

\section{Speaker's guide for the 90-minute talk}

Spend the first 10 minutes making the universal viewpoint clear.  The audience should understand that computing $H^*(\Gn;\F_2)$ is the same as computing all mod $2$ characteristic classes of rank $n$ real bundles.  Emphasize that the talk assumes existence of Stiefel--Whitney classes and proves the ring computation plus uniqueness, not existence.

From 10 to 25 minutes, list the inputs without proving them: classification of vector bundles by $\Gn$, the computation $H^*(\RP^\infty;\F_2)=\F_2[a]$, the Kunneth theorem for $(\RP^\infty)^n$, the Schubert cell count, and the fundamental theorem of symmetric polynomials.  Do not let this section expand.

From 25 to 45 minutes, focus on the split bundle
\[
  \xi=\ell_1\oplus\cdots\oplus\ell_n
\]
over $(\RP^\infty)^n$.  This is the conceptual heart of the first half.  Say out loud: ``We are testing the universal bundle on the most split possible rank $n$ bundle.''  Then compute
\[
  w(\xi)=\prod_i(1+a_i)
\]
and conclude that
\[
  g^*w_i=e_i(a_1,
  \dots,a_n).
\]
When proving injectivity, stress that this proves algebraic independence of the $w_i$'s, not yet injectivity of $g^*$ on all cohomology.

From 45 to 60 minutes, present the dimension comparison slowly.  Put the two counts side by side:
\[
  \text{Schubert cells in degree $r$}
  \longleftrightarrow
  \text{partitions of $r$ into at most $n$ parts},
\]
while
\[
  \text{degree $r$ monomials in $w_1,\dots,w_n$}
  \longleftrightarrow
  \text{partitions of $r$ with largest part at most $n$}.
\]
Then draw one Young diagram and its conjugate.  This is worth the board time.

From 60 to 68 minutes, do only one or two low-rank examples.  The case $n=2$ is probably enough:
\[
  H^*(G_2;\F_2)=\F_2[w_1,w_2].
\]
Use it to make the phrase ``no relations'' concrete.

From 68 to 82 minutes, prove uniqueness.  Use the notation $W$ and $\olW$.  The audience must hear the four steps clearly: line bundles, split test bundle, universal bundle, arbitrary pullback.  The most important sentence is: ``After the ring computation, $g^*$ is injective because every class is a polynomial in the $w_i$'s and $g^*w_i=e_i$.''

Reserve the last 8 minutes for the final summary and questions.  A good closing sentence is: ``The split test bundle detects the generators, and the Schubert cell count proves there are no hidden classes.''

\section{Final summary}

The main theorem is
\[
  H^*(\Gr_n(\R^\infty);\F_2)
  \cong
  \F_2[w_1,\dots,w_n].
\]
Here $w_i=w_i(\gamman)$ is the $i$-th Stiefel--Whitney class of the tautological rank $n$ bundle.

The proof has two key ideas.  The first is the splitting test over $(\RP^\infty)^n$, where the total Stiefel--Whitney class becomes
\[
  \prod_{i=1}^n(1+a_i).
\]
This shows that the universal classes pull back to elementary symmetric polynomials and therefore satisfy no algebraic relation.  The second is Schubert cell counting: in each degree, the number of monomials in the $w_i$'s equals the number of Schubert cells.  So the polynomial subring generated by the $w_i$'s is already the whole cohomology ring.

As an application, the same computation proves uniqueness of Stiefel--Whitney classes.  Any two assignments satisfying the axioms agree on line bundles, hence on split bundles over $(\RP^\infty)^n$, hence on the universal bundle $\gamman$, and finally on every bundle by pullback along a classifying map.

\begin{thebibliography}{9}

\bibitem{MilnorStasheff}
John Milnor and James D. Stasheff.
\textit{Characteristic Classes}.
Annals of Mathematics Studies, No. 76. Princeton University Press, 1974.

\bibitem{HatcherVBKT}
Allen Hatcher.
\textit{Vector Bundles and K-Theory}.
2017. Available from the author's website.

\end{thebibliography}

\end{document}
